% 第3節 命題と証明
\section{命題と証明}

\subsection{命題の逆, 裏, 対偶}

% 必要に応じて内容を追加してください

\newpage

\subsection{対偶を利用する証明}

% 必要に応じて内容を追加してください

\newpage

\subsection{背理法を利用する証明}

ある命題について, その命題が成り立たないと仮定して, 矛盾を導くことでもとの命題を証明する方法がある.

% 例題2
\begin{example}[【教p.67 例題2】]
    $\sqrt{2}$ が無理数である $^{*}$ことを用いて, 次の命題を証明せよ.
    \[1 + \sqrt{2} \text{ は無理数である}\]
\end{example}

\begin{proof}\mbox{}\\\hspace{1em}%
    $1 + \sqrt{2}$ が無理数でないと仮定すると, $1 + \sqrt{2}$ は有理数である.

    その有理数を $r$ とすると, $1 + \sqrt{2} = r$ より,
    \[\sqrt{2} = r - 1\]
    $r$ が有理数ならば, $r - 1$ も有理数であるからこの等式は $\sqrt{2}$ が無理数であることに矛盾する.

    したがって, $1 + \sqrt{2}$ は無理数である.
\end{proof}
\vspace{3em}

例題の証明は, 次のような流れである.
\begin{enumerate}
    \item \label{enum:cond1} 命題が成り立たないと仮定する.
    \item \label{enum:cond2}\ref{enum:cond1} の仮定のもとで矛盾を導く.
    \item \ref{enum:cond2} で矛盾が生じたのは, \ref{enum:cond1} の仮定が間違っているからである.
    \item したがって, もとの命題が成り立つ.
\end{enumerate}

このような証明方法を \textbf{背理法} という.
\vspace{1em}

\begin{theorembox}[thm:背理法を利用する証明]{\textbf{背理法を利用する証明}}
    命題が成り立たないと仮定して矛盾を導くことにより、もとの命題が真であると結論する。
\end{theorembox}

\vfill

\noindent
\rule{\textwidth}{0.4pt}

\noindent
$^{*}$ $\sqrt{2}$ が無理数であることは, 後で証明することとする.
\newpage

% 練習19
\begin{exercise}[【教p.67 練習19】]
    $\sqrt{2}$ が無理数であることを用いて, 次の命題を証明せよ.
    \[1 + 3\sqrt{2} \text{ は無理数である}\]
\end{exercise}

\begin{problemsolution}\\\hspace{1em}%
    $1 + 3\sqrt{2}$ が無理数でないと仮定すると, $1 + 3\sqrt{2}$ は有理数である.

    その有理数を $r$ とすると, $1 + 3\sqrt{2} = r$ より,

    \begin{align*}
        3\sqrt{2} &= r - 1 \\
        \sqrt{2} &= \dfrac{r - 1}{3}
    \end{align*}
    $r$ が有理数ならば, $\dfrac{r - 1}{3}$ も有理数であるからこの等式は $\sqrt{2}$ が無理数であることに矛盾する.

    したがって, $1 + 3\sqrt{2}$ は無理数である.
\end{problemsolution}
\newpage

\subsection{研究 \quad $\sqrt{2}$ が無理数であることの証明}

前ページの例題2では、「$\sqrt{2}$ が無理数である」ということを認めていた。

この事実を背理法を利用して証明してみよう。

\begin{proof}\mbox{}\\\hspace{1em}%
    「$\sqrt{2}$ が無理数でない」すなわち「$\sqrt{2}$ は有理数である」と仮定すると,

    $\sqrt{2}$ はあり自然数 $m , n$ を用いて

\begin{equation}\tag{1}\label{eq:1}
    \sqrt{2} = \dfrac{m}{n}
\end{equation}
\vspace{1em}

    と表すことができる.このとき、できる限り約分して、$m$ と $n$ に1以外の正の約数がないような分数$^{*}$にする。

    \eqref{eq:1} から,

    \begin{align*}
        \sqrt{2} \; n = m
    \end{align*}
    \vspace{1em}

    この両辺を2乗すると,

    \begin{equation}\tag{2}\label{eq:2}
        2n^2 = m^2
    \end{equation}
    \vspace{1em}

    よって、$m^2$ は偶数である。

    $m^2$ が偶数ならば, \uwave{$m$ も偶数} である。     %p66の例題1より〜

    偶数 $m$ はある自然数 $k$ を用いて、 $m = 2k$ と表されるから、\ref{eq:2} に代入して,

    \begin{align*}
        2n^2 &= 4k^2\\
        \intertext{すなわち}
        n^2 &= 2k^2
    \end{align*}
    \vspace{1em}

    よって、$n^2$ は偶数となり、\uwave{$n$ も偶数} である。

    $m$ と $n$ がともに偶数となることは, $m$ と $n$ に1以外の正の約数がないとしたことに矛盾する.

    したがって, $\sqrt{2}$ は無理数である.
\end{proof}


\vfill

\noindent
\rule{\textwidth}{0.4pt}

\noindent
$^{*}$ このような分数を既約分数という。
\newpage
